\documentclass[10pt,a4paper]{article}
\usepackage{amsmath,amssymb,amsthm,graphicx}
% Minimal retained preamble for S3 SIGMA article cards taken from
% sources/sigma-2023-026-on-generalized-wkb-expansion-of-monodromy-ge/pinned-original-source.tex
% (doi:10.3842/SIGMA.2023.026, SIGMA 19 (2023), CC BY 4.0).
%
% Omitted and why: the publisher \documentclass{sigma} block and the
% frontmatter commands (FirstPageHeading, ShortArticleName, ArticleName,
% Author, AuthorNameForHeading, Address, Email, ArticleDates, Abstract,
% Keywords, Classification) - the sigma class is unavailable in the
% toolchain and the assembler supplies the article class; the tikz-cd
% package is dropped because no counted or context span of this card
% uses a tikzcd diagram (unused-package disclosure per the three-pitfall
% broadcast, item ii).
% The theorem declarations, equation numbering and paper macros below
% are copied verbatim from the pinned source (source lines 9-38).
% hyperref is added beyond the source preamble because the author's
% bibliography entries use \href and the unavailable sigma class is
% what provides it in the published version.
\usepackage{amsmath,amssymb,amsthm}
\usepackage{mathtools}
\usepackage{hyperref}
\providecommand{\qed}{\ensuremath{\square}}
\allowdisplaybreaks
\numberwithin{equation}{section}
\newtheorem{Theorem}{Theorem}[section]
\newtheorem{Lemma}[Theorem]{Lemma}
\newtheorem{Proposition}[Theorem]{Proposition}
\newtheorem{Corollary}[Theorem]{Corollary}

\theoremstyle{definition}
\newtheorem{Definition}[Theorem]{Definition}
\newtheorem{Remark}[Theorem]{Remark}

\makeatletter
\newcommand{\hathat}[1]{%
\begingroup%
 \let\macc@kerna\z@%
 \let\macc@kernb\z@%
 \let\macc@nucleus\@empty%
 \hat{\raisebox{.2ex}{\vphantom{\ensuremath{#1}}}\smash{\hat{#1}}}%
\endgroup%
}
\makeatother
%-------------------------------------------------
\newcommand{\R}{\mathbb{R}}
\newcommand{\N}{\mathbb{N}}
\newcommand{\Z}{\mathbb{Z}}
\newcommand{\CC}{\mathbb{C}}
\newcommand{\T}{\mathbb{T}}
\title{Homological 2-form equivalence and the generating function of the period symplectomorphism (Proposition 3.4)}
\author{Klimov, Roman}
\date{doi:10.3842/SIGMA.2023.026; CC BY 4.0}
\begin{document}
\maketitle
\setcounter{section}{3}
\setcounter{Theorem}{3}
\noindent\textit{Editorial note.} All delivered mathematical text is the author's own, copied from the pinned source file sources/sigma-2023-026-on-generalized-wkb-expansion-of-monodromy-ge/pinned-original-source.tex; no statement or proof step was rewritten, reordered or removed.\par\medskip
\noindent\textit{Editorial note.} The counted span is the author Proposition environment Prop\_3 with its own label and its own single proof environment (Proof of Proposition 3.4, source lines 661-832).\par\medskip
\noindent\textit{Editorial note.} The statement applies the Theta construction to the family S\_0 - S\_1; the construction from source lines 131-154 (printed (12) and (form)) is delivered as a context span before the statement. The regularization (76), the pairing (77), the variational rule (80) with the coordinate (36) it cites, the involutive cycles (32)/(33), the pole coordinate (39) and the variational formulas (51)-(54) are cited by the proof and delivered as context spans after the statement, so they print locally renumbered (Remark 3.5, Lemma 3.6, Remark 3.7, Proposition 3.8) while the statement keeps 3.4.\par\medskip
\noindent\textit{Editorial note.} The proof cites Lemma 3.6 (source line 740); the Lemma 3.6 statement is delivered as a context span so \textbackslash{}ref\{Lem\_1\} resolves (it prints locally renumbered, keeping the printed number 3.6 coincidentally), and its full proof is produced as the sibling card of this batch. The Riemann bilinear identity, Stokes' theorem and the Poincare lemma are named standard facts; connectedness of the stratum is the cited Chen-Gendron Theorem 1.3 (bibliography fragment). These are disclosed as named prerequisites.\par\medskip
\noindent\textit{Editorial note.} The bibliography fragment bib.tex carries the two cited bibitems (korotkin2018periods, chen2021classification) extracted verbatim from the pinned source's own thebibliography (source lines 1471-1475 and 1404-1408); hyperref is added to the card preamble to provide \textbackslash{}href.\par\medskip
\noindent\textit{Editorial note.} Printed number anchored to the published PDF (SIGMA 19 (2023) 026, 36 pages, read from the stream-extracted text of sigma23-026.pdf downloaded from sigma-journal.com). The pinned tex has four \textbackslash{}section commands and no \textbackslash{}appendix and its section structure matches the published PDF one-to-one, so the unit 010 pinned-versus-published realignment does not trigger.\par\medskip
\noindent\textit{Editorial note.} Equation and theorem numbers inside the card are locally renumbered; every cross-reference is a symbolic \textbackslash{}eqref/\textbackslash{}ref and resolves inside the card.\par\medskip
\noindent\textit{Editorial scope.} Single environment card cut from the SIGMA article doi:10.3842/SIGMA.2023.026 (sigma-2023-026-on-generalized-wkb-expansion-of-monodromy-ge); statement and proof are the author's own text at source lines 510-529 and 661-832.\par\medskip
\noindent\textit{Allowed standard prerequisites (editorial).}\begin{itemize}
\item Riemann bilinear identity. Standard named fact: for meromorphic differentials on a compact Riemann surface, the period pairing equals the sum of residues of the product against a primitive inside a fundamental polygon.. Load-bearing: the proof rewrites the pairing of (oint v\_0, oint Qtilde/v\_0) as a residue sum (109) and reduces the hbar\textasciicircum{}k coefficients to residues (121) and (122).
\item Variational formulas (51), (52). Delivered in context (author Proposition environment, source lines 398-404): derivatives of the period matrix and of normalized holomorphic differentials against period coordinates are contour integrals.. Load-bearing: the proof evaluates the pairing of Qtilde/v\_0 against delta oint v\_0 via (111) and (114).
\item Period-pairing variational identity (Lemma 3.6). Delivered in the unit as its own card (Lem\_1): the identity (78)/(79) rewriting the pairing with delta oint w\_2 as a boundary integral plus a (w\_1 w\_2/v)-pairing term.. Load-bearing: the proof cites Lemma 3.6 (source line 740) to rewrite the pairings in (106) as (116) and (117).
\item Connectedness of the stratum. Theorem 1.3 of Chen-Gendron (cited as chen2021classification): generically (outside the hyperelliptic locus for g >= 3) the stratum Q\_\{g,n\}, hence Q\_\{g,n\}[r], is connected.. Load-bearing: the arbitrary-pair case joins Q\_0 and Q\_1 by a path and applies the nearby-pair identity telescopically (source lines 816-829).
\item Stokes' theorem and the Poincare lemma. Standard named facts of multivariable calculus on manifolds.. Load-bearing: (79) follows from (78) by Stokes (source lines 654-657 of the Lemma 3.6 proof) and the closed-form hypothesis is integrated locally by the Poincare lemma (source lines 831-832).
\end{itemize}\medskip
\begin{equation}\label{66}Q_1-Q_0=\frac{1}{2}(S_0 -S_1 ).\end{equation}

Introduce the Bergman projective connection $S_B$ defined in terms of the canonical bidifferential $B(x,y)$ on $\mathcal{C}$, which is normalized with respect to chosen Torelli marking in $H_1(\mathcal{C},\Z)$:%
\begin{equation}\label{12}
 S_B(x)=\left(B(x,y)-\frac{{\rm d}\xi(x) {\rm d}\xi(y)}{(\xi(x)-\xi(y))^2}\right)\Big|_{y=x},
\end{equation}
where $\xi$ is any local coordinate near point $x.$ Since $S_B$ depends holomorphically on the conformal structure of $\mathcal{C},$ the difference $S-S_B$ becomes
a family of quadratic differentials with at most simple poles at the punctures $(z_j)$, depending holomorphically on moduli of $\mathcal{M}_{g,n}.$ Using the identification of the moduli space of quadratic differentials with simple poles and the cotangent bundle $T^*\mathcal{M}_{g,n}$, we can associate $S-S_B$ with the 1-form $\Theta_{(S-S_B)}$, locally defined on $\mathcal{M}_{g,n}$, in the following way.

At first, introduce the set of holomorphic local coordinates $(\Omega_{jk}, q_l)$ on $\mathcal{M}_{g,n}, \, g \geq 2$. To determine locally the
conformal structure of $\mathcal{C}$ we pick at generic point of $\mathcal{M}_{g,n}$ (outside of hyperelliptic locus for $g \geq 3$) a set $D$ of $3g-3$ entries of the
period matrix $\Omega$ of $\mathcal{C}.$ The quadratic differentials corresponding to cotangent vectors $\delta \Omega_{jk} $ are products $u_j u_k$ of normalized
holomorphic differentials. An additional set of $n$ coordinates which determine the positions of punctures $(z_l)^n_{l=1} $
on $\mathcal{C}$ we choose to be $q_l = (u_i/u_j )(z_l)$ where $u_i$ and $u_j$ form a pair of normalized holomorphic 1-forms on
$\mathcal{C}$, such that $u_j(z_l) \neq 0$. The quadratic differential corresponding to cotangent vector $\delta q_l$ is the meromorphic quadratic
differential $Q^{z_l}$ (given by the formula \eqref{110} below) whose only simple pole is at $z_l$. These coordinates are local: in different coordinate charts on
$\mathcal{M}_{g,n}$ one might need to choose other pairs of normalized holomorphic differentials and/or different Torelli
markings. The momenta $p_l$ are then defined
to be coefficients of decomposition of the quadratic differential $S-S_B$ in the
basis described above.
Writing down the quadratic differential $S-S_B$ as
\begin{equation*}
S-S_B=\sum_{(jk) \in D}p_{jk} u_j u_k+ \sum^n_{l=1} p_l Q^{z_l},
\end{equation*}
where $p_{jk}$ and $p_l$ are holomorphic functions of $(\Omega_{jk},q_l)$, the corresponding 1-form $\Theta_{(S-S_B)}$ on~$\mathcal{M}_{g,n}$ reads as
\begin{equation}\label{form}\Theta_{(S-S_B)}=\sum_{(jk) \in D}p_{jk} \delta \Omega_{jk}+ \sum^n_{l=1} p_l \delta q_l.\end{equation}

\begin{Proposition} \label{Prop_3}\quad
\begin{enumerate}\itemsep=0pt
\item[$1.$] Two meromorphic differentials $Q_0, Q_1 \in \mathcal{Q}_{n,g}[\textbf{r}]$ induce the same homological $2$-form on $\mathcal{Q}_{n,g}[\textbf{r}]$ if and only if the $1$-form $\Theta_{(S_0-S_1)}$, corresponding to family of quadratic differentials $S_0-S_1$ and locally defined on $\mathcal{M}_{g,n}$,
is closed, $\delta\Theta_{(S_0-S_1)}=0.$

\item[$2.$] The generating function of the symplectomorphism between the periods $\big(A^{(0)}_k,B^{(0)}_k\big)$ of $v_0$ and $\big(A^{(1)}_k,B^{(1)}_k\big)$ of $v_1$ for the chosen potentials \begin{equation}\label{69}\theta_0=\sum^{g^-}_{k=1}\big(B^{(0)}_k \delta A^{(0)}_k - A^{(0)}_k \delta B^{(0)}_k\big), \qquad \theta_1=\sum^{g^-}_{k=1}\big(B^{(1)}_k \delta A^{(1)}_k - A^{(1)}_k \delta B^{(1)}_k\big),\end{equation}
defined by
\begin{equation*}%\label{70}
\delta\mathcal{G}_{{\rm hom}}=\mathcal{H}^*\theta_1 - \theta_0 ,
\end{equation*}
has the following form:
\begin{equation}\label{71}
\mathcal{G}_{{\rm hom}}=\sum^n_{i=1}\pi {\rm i} r_i\left({\rm reg}\int^{z^{(1)}_i}_{z^{(2)}_i}v_1 -{\rm reg}\int^{z^{(1)}_i}_{z^{(2)}_i}v_0 \right)+\frac{1}{2}\Hat{G}_{(S_0 -S_1)},
\end{equation}
where there exists a local holomorphic function $\Hat{G}_{(S_0 -S_1)}$ on $\mathcal{M}_{g,n}$, such that
\begin{equation*}%\label{72}
\delta\Hat{G}_{(S_0 -S_1)}= \Theta_{(S_0-S_1)}.
\end{equation*}
\end{enumerate}
\end{Proposition}

\begin{Remark}
For any $Q \in \mathcal{Q}_{g,n}[\mathbf{r}]$ the integral $\int^{z^{(1)}_i}_{z^{(2)}_i} v$ is singular at the endpoints. We define its regularization by removing the divergent part as follows: fix a coordinate $\xi_j$ near $z_j$ such that
\begin{equation*}%\label{73}
Q(x) \sim \left(\frac{r_j^2}{\xi_j^2} +O\big(\xi_j^{-1}\big)\right)({\rm d} \xi_j)^2.\end{equation*}
$\xi_j$ can also serve as a local coordinate on $\hat{\mathcal{C}}$ near the lifts $\big\{ {z^{(1)}_j}, {z^{(2)}_j}\big\}$ with
\begin{equation*}%\label{74}
v(x)\sim \pm \left(\frac{ r_j}{\xi_j} +O(1)\right){\rm d} \xi_j.\end{equation*}
Let $z^t_j$ be an arbitrary sequence of points on $\mathcal{C}$ converging to $z_j,$ such that in the local coordinate~$\xi_j$
\begin{equation*}%\label{75}
\operatorname{Re} \big(\xi_j\big(z^t_j\big)\big) \sim \frac{1}{t}, \quad t \xrightarrow{} \infty, \qquad \operatorname{Im} \big(\xi_j\big(z^t_j\big)\big)=0.\end{equation*}
Then the regularization is defined by
\begin{equation}\label{76}
{\rm reg}\int^{z^{(1)}_j}_{z^{(2)}_j}v := \lim_{t \xrightarrow{} \infty } \left( \int^{z^{t (1)}_j}_{z^{t (2)}_j}v - 2 r_j \log t \right).
\end{equation}
\end{Remark}

\begin{equation}\label{77} \left\langle\oint w_1,\oint w_2 \right\rangle:=\sum^{g^-}_{j=1}\left[\oint_{b^-_j}w_1 \oint_{a^-_j}w_2 - \oint_{a^-_j}w_1 \oint_{b^-_j}w_2\right].\end{equation}

\begin{Lemma} \label{Lem_1}
Let $w_1$ and $w_2$ be two meromorphic differentials on $\hat{\mathcal{C}}$, skew-symmetric under involution. Assuming both $w_1$ and $w_2$ holomorphically depend on moduli $(A_i, B_i)^{g^-}_{i=1}$, the following identity of $1$-forms on $\mathcal{Q}_{g,n}[\textbf{r}]$ holds:
\begin{equation}\label{78}\left\langle \oint w_1, \delta \oint w_2\right\rangle =-\frac{1}{2}\int_{\partial \hat{\mathcal{C}}_0}\left(\delta w_2
\int^x_{p_0} w_1\right)+\left\langle \oint \frac{w_1 w_2}{v},\delta \oint v\right\rangle \end{equation} or
\begin{equation}\label{79}\left\langle \oint w_1, \delta \oint w_2\right\rangle =\frac{1}{2}\int_{\partial \hat{\mathcal{C}}_0}\left( w_1
\int^x_{p_0} \delta w_2\right)+\left\langle \oint\frac{w_1 w_2}{v},\delta \oint v\right\rangle. \end{equation}
Here $\hat{\mathcal{C}_0}$ is the fundamental polygon of the covering surface $\hat{\mathcal{C}}$, $p_0$ is a generic point.
\end{Lemma}

\begin{Remark}
Note that by $\delta \oint w$ appearing in the pairing we mean differential applied to the periods of $w$, while $\delta w$ is given by
\begin{equation}\label{80}\delta w:=\sum^{g^-}_{i=1} \left(\frac{\delta w}{\delta A_i}\Big|_{z(x)} \delta A_i +\frac{\delta w}{\delta B_i}\Big|_{z(x)} \delta B_i\right),\end{equation}
where the partial derivatives are defined under the assumption that the local coordinate~\eqref{36} is independent of the moduli. For example,
\begin{equation*}\frac{\delta w}{\delta A_i}\Big|_{z(x)}:=
v(x) \frac{\delta}{\delta A_i}\Big|_{z(x)={\rm const}}\left\{\frac{w(x)}{v(x)}\right\}, \end{equation*}
where $w(x)/v(x)$ is
a meromorphic function on $\hat{\mathcal{C}}$.
\end{Remark}

\begin{equation}\label{36}
z(x)=\int^x_{x_1}v,
\end{equation}

\begin{gather}\label{32}a^-_k=\frac{1}{2}\big(a_k-a^\mu_k\big), \qquad b^-_k=\frac{1}{2}\big(b_k-b^\mu_k\big),\\
\label{33}a^-_l=\frac{1}{\sqrt{2}}\tilde{a}_l, \qquad b^-_l=\frac{1}{\sqrt{2}}\tilde{b}_l,\\
%\label{34}
t^-_j=\frac{1}{2}\big(t_j-t^\mu_j\big)\nonumber
\end{gather}

\begin{equation}\label{39}\zeta_i(x)=\exp \left(\frac{1}{ r_i} \int^x_{x_1}v\right). \end{equation}

\begin{Proposition}
For arbitrary basis $\{s_j\}^{6g-6+2n}_{j=1}$ of $H_-$ and its dual basis $\{ s^{*}_{j} \}^{6g-6+2n}_{j=1}$ the following formulas hold on $\mathcal{Q}_{g,n}[\mathbf{r}]$:
\begin{gather}\label{51}\frac{\delta \Omega_{ij}}{\delta \mathcal{P}_{s}}=\frac{1}{2}\oint_{s^*}{f}_i {f}_j v,\\
\label{52}\frac{\delta f_j(x)} {\delta \mathcal{P}_s}\Big|_{z(x)={\rm const}}=\frac{1}{4 \pi {\rm i}}\oint_{s^*}{f}_j(t){b}(x,t)v(t),\\
\label{53}\frac{\delta b(x,y)}{\delta \mathcal{P}_s}\Big|_{z(x),z(y)={\rm const}}=\frac{1}{4 \pi {\rm i}}\oint_{s^*}{b}(x,t){b}(t,y)v(t),\\
\label{54}\frac{\delta q(x)}{\delta \mathcal{P}_s}\Big|_{z(x)={\rm const}}=\frac{3}{4 \pi {\rm i}}\oint_{s^*}{h}(x,t)v(t).\end{gather}
\end{Proposition}

\begin{proof}[Proof of Proposition \ref{Prop_3}]
\textit{Case of two nearby differentials} $Q_0, Q_1 \in \mathcal{Q}_{g,n}[\mathbf{r}] $: let $U$ be a~sim\-ply-connected neighborhood of $Q_0$ and take $Q_1 \in U.$ For sufficiently small $\hbar$ this differential may be expressed as $Q_1=Q_0+\hbar \tilde{Q},$ where $\tilde{Q}$ is a quadratic differential with at most simple poles at~$(z_j)^{n}_{j=1}$. The canonical cover $\hat{\mathcal{C}_\hbar}$, defined by $\big(v^\hbar_1\big)^2=Q_1$, becomes $\hbar$-dependent. Consider $\mathcal{P}_{\hbar}$ to be one of the periods $\big(A^{(1)}_i,B^{(1)}_i\big)$ of $v^\hbar_1$. Then its $k$-th derivative with respect to $\hbar$ is given~by\footnote{Double factorial is defined by $n!!:=n(n-2)(n-4)\cdots$.}
\begin{equation}\label{101}
\frac{\partial^k }{\partial \hbar^k}\mathcal{P}_{\hbar}\Big|_{\hbar=0} =
(-1)^{k+1}\frac{(2k-3)!!}{2^k}\oint_{s}\frac{\tilde{Q}^k}{v_{0}^{2k-1}},\end{equation}
where $s$ is an element of $H_{-}(\hat{\mathcal{C}}_0, \Z)$.

To justify these formulas consider $ \mathcal{P}_{\hbar}=\oint_{s(\hbar)}v^\hbar_{1}$ as the integral on the base curve $\mathcal{C}$ via the projection $\pi\colon \hat{\mathcal{C}}_\hbar \xrightarrow{} {\mathcal{C}}$. If cycle $s(\hbar)$ belongs to the subset~\eqref{32}, using skew-symmetry of $v^\hbar_1$ it projects onto one of homology cycles of $\mathcal{C}$, which is independent of $\hbar$. If $s(\hbar)$ is an element of~\eqref{33}, it projects onto the contour encircling or passing through branch cuts arranged between pairs of zeroes of $Q^\hbar_{1}$. Despite the positions of the branch points vary along with $\hbar$, we may assume that the projections $\pi(s(\hbar))$ are kept fixed on $\mathcal{C}$. Then one has
\begin{equation*}%\label{102}
\frac{\partial}{\partial \hbar}\oint_{s(\hbar)}v^\hbar_{1}=\frac{\partial}{\partial \hbar}\oint_{\pi(s(\hbar))}\sqrt{Q_0+\hbar \tilde{Q}}= \oint_{\pi(s(\hbar))}\frac{\partial}{\partial \hbar}\sqrt{Q_0+\hbar\tilde{Q}},
\end{equation*}
and the differentiation is followed by pullback to $\hat{\mathcal{C}}$. Higher derivatives are obtained the same way.

Applying this argument, we can expand period coordinates by powers of $\hbar.$ Write
\begin{align*}%\label{103}
\mathcal{P}_{\hbar}& = \mathcal{P}_{\hbar}\Big|_{\hbar=0}+ \frac{\partial}{\partial \hbar}\mathcal{P}_{\hbar}\Big|_{\hbar=0}\hbar+\dots +\frac{\partial^k}{\partial \hbar^k}\mathcal{P}_{\hbar}\Big|_{\hbar=0}\frac{\hbar^k}{k!}+\cdots\\
%\label{104}
& =\oint_{s} {v_0} + \frac{\hbar}{2}\oint_{s} \frac{\tilde{Q}}{ v_0}+\dots +\hbar^k\binom{\frac{1}{2}}{k}\oint_{s}\frac{\tilde{Q}^k}{v_0^{2k-1}}+\cdots.
\end{align*}
Plugging $\hbar$-expansions of the periods $\big(A^{(1)}_i,B^{(1)}_i\big)$ into the potential $\theta_1$ \eqref{69} and arranging terms by powers of $\hbar$, we write with the help of the pairing notation \eqref{77}:
\begin{gather}\label{105}
\theta_1=\left\langle \oint\! v_{1}, \delta \oint\! v_{1}\right\rangle =\left\langle \oint\! v_0, \delta \oint\! v_0\right\rangle + \hbar \left[\frac{1}{2}\left\langle \oint \!v_0, \delta \oint\! \frac{\tilde{Q}}{v_0}\right\rangle +\frac{1}{2}\left\langle \oint\! \frac{\tilde{Q}}{v_0}, \delta \oint\! v_0\right\rangle \right]\!\!\!\\
\hphantom{\theta_1=}{}
+\sum^{\infty}_{k=2}\hbar^k\Bigg[\binom{\frac{1}{2}}{k}\left(\left\langle \oint v_0,\delta \oint \frac{\tilde{Q}^k}{v_0^{2k-1}}\right\rangle +\left\langle \oint \frac{\tilde{Q}^k}{v_0^{2k-1}},\delta \oint v_0\right\rangle \right)\nonumber
\\
\hphantom{\theta_1=}{}
+\sum^{k-1}_{l=1} \binom{\frac{1}{2}}{l} \binom{\frac{1}{2}}{k-l}
\left\langle \oint \frac{\tilde{Q}^l}{v_0^{2l-1}}, \delta \oint \frac{\tilde{Q}^{k-l}}{v_0^{2(k-l)-1}} \right\rangle \Bigg].\label{106}
\end{gather}

We will treat separately the expressions near $\hbar^1$ and $\hbar^k$, $k \geq 2$.

\textit{Coefficient near $\hbar^1$}: noticing that
\begin{equation*}%\label{107}
\delta\left\langle \oint v_0, \oint \frac{\tilde{Q}}{v_0}\right\rangle =\left\langle \delta \oint v_0, \oint \frac{\tilde{Q}}{v_0}\right\rangle +\left\langle \oint v_0, \delta \oint \frac{\tilde{Q}}{v_0}\right\rangle,
\end{equation*}
the expression near $\hbar^1$ could be rewritten as
\begin{equation*}%\label{108}
\frac{1}{2}\delta\left\langle \oint v_0, \oint \frac{\tilde{Q}}{v_0}\right\rangle +\left\langle \oint \frac{\tilde{Q}}{v_0}, \delta \oint v_0\right\rangle .
\end{equation*}
 Applying the Riemann bilinear identity, the pairing $\big\langle \oint v_0, \oint \frac{\tilde{Q}}{v_0}\big\rangle $ could be written as the sum over residues inside the fundamental domain. Differential $\tilde{Q}$, being lifted to $\hat{\mathcal{C}}$, gains double zeroes at branch points $x_i$, which are the zeroes of $Q_0$, and simple poles at preimages of~$z_j$. That makes $\frac{\tilde{Q}}{v_0}$ holomorphic, while $v_0$ has poles at
$\big\{z^{(1)}_j, z^{(2)}_j \big\}^n_{j=1}$. Therefore, using that near $z^{(1)}_j$ \big($z^{(2)}_j$\big) in the local coordinate~\eqref{39} $v_0=\pm \frac{r_j}{\zeta_i}{\rm d}\zeta_i$, we write
\begin{equation}\label{109}\delta
\left[\sum^n_{j=1}\pi {\rm i} \underset{\{z^{(1)}_j, z^{(2)}_j \}}{\rm res}\left(v_0
\int^x_{p_0}\frac{\tilde{Q}}{v_0} \right)\right]=\delta\left[\sum^n_{j=1}\pi{\rm i} r_j
\int^{z^{(1)}_j}_{z^{(2)}_j}\frac{\tilde{Q}}{v_0} \right]. \end{equation}
We will also express the pairing $\big\langle \oint \frac{\tilde{Q}}{v_0}, \delta \oint v_0\big\rangle $ in a different form, introducing the system of local coordinates on $\mathcal{M}_{g,n}$. For simplicity here we restrict us to the case $g \geq 2$ (low genus cases $g=0,1$ could be covered by analogy following~\cite{korotkin2018periods}). At generic point of the moduli space~$\mathcal{M}_{g,n}$ the quadratic differential~$\tilde{Q}$ could be represented as a linear combination of $3g-3$ products of normalized holomorphic differentials $u_j u_k$, where $(jk) \in D$ for some subset~$D$ of entries of matrix~$\Omega$, and additional~$n$ quadratic differentials encoding the meromorphic part could be represented by the following generically meromorphic differentials~$Q_k$ whose only pole of order one located at~$z_k$:
\begin{equation}\label{110}
Q^{z_k}(t)=\frac{1}{4 \pi {\rm i}}\frac{u_i(t)u_j(z_k)-u_i(z_k)u_j(t)}{u_j^2(z_k)}B(t, z_k),
\end{equation}
here $u_i$ and $u_j$ are two arbitrary normalized holomorphic differentials such that $u_j(z_k) \neq 0$.

Using the variational formulas \eqref{52} after lifting the function $\frac{u_i}{u_j}(x)$ to $\hat{\mathcal{C}}$, we get
\begin{equation}\label{111}
\frac{\delta }{\delta \mathcal{P}_s}\left[\frac{u_i}{u_j}(z_k)\right]=\frac{\delta }{\delta \mathcal{P}_s}\Bigg[\frac{u_i}{u_j}\big(z^{(1)}_k\big)\Bigg]=\oint_{s^*}\frac{Q^{z^{(1)}_k}}{v_0},
\end{equation}
where $Q^{z^{(1)}_k}$ is a lift of $Q^{z_k}$ to $\hat{\mathcal{C}}$. The entries $\Omega_{jk}, \ (jk) \in D$ of the period matrix can serve as the moduli of the base curve $\mathcal{C},$ while $\frac{u_i}{u_j}(z_k):=q_k$ code the positions of poles, providing in total $3g-3+n$ local coordinates on $\mathcal{M}_{g,n}.$

At generic point of $\mathcal{M}_{g,n}$ quadratic differential $\tilde{Q}$ can be expressed as
\begin{equation*}%\label{112}
\tilde{Q}=\sum_{(jk) \in D}p_{jk} u_j u_k+ \sum^n_{l=1} p_l Q^{z_l}, \qquad p_{jk}, p_l \in \CC.\end{equation*}
Then applying variational formulas \eqref{51} and \eqref{111}, one has
\begin{gather}
\left\langle \oint \frac{\tilde{Q}}{v_0}, \delta \oint v_0\right\rangle =\sum^{g_-}_{j=1} \left[ \left(\oint_{b^-_j}\frac{\tilde{Q}}{v_0} \right)\delta A^{(0)}_j - \left(\oint_{a^-_j}\frac{\tilde{Q}}{v_0} \right)\delta B^{(0)}_j \right] \nonumber\\
\qquad{}
 =\sum_{(jk) \in D}p_{jk}\sum^{g_-}_{j=1} \left[ \left(\oint_{b^-_j}\frac{u_j u_k}{v_0} \right)
 \delta A^{(0)}_j - \left(\oint_{a^-_j}\frac{u_j u_k}{v_0} \right)\delta B^{(0)}_j \right ]\nonumber\\
 \qquad\quad{} + \sum^n_{l=1} p_l \sum^{g_-}_{j=1} \left[ \left(\oint_{b^-_j}\frac{Q^{z^{(1)}_l}}{v_0} \right)\delta A^{(0)}_j
 - \left(\oint_{a^-_j}\frac{Q^{z^{(1)}_l}}{v_0} \right)\delta B^{(0)}_j \right ]\nonumber\\
\qquad{}
=\sum_{(jk) \in D}\!p_{jk}\sum^{g_-}_{j=1} \!\left[ \frac{\partial \Omega_{ij}}{\partial A^{(0)}_j} \delta A^{(0)}_j +\frac{\partial \Omega_{ij}}{\partial B^{(0)}_j}\delta B^{(0)}_j \right]+\sum^n_{l=1} p_l \sum^{g_-}_{j=1} \!\left[ \frac{\partial q_l}{\partial A^{(0)}_j}\delta A^{(0)}_j + \frac{\partial q_l}{\partial B^{(0)}_j}\delta B^{(0)}_j \right ]
\nonumber\\
\qquad{}
=\sum_{(jk) \in D}p_{jk} \delta \Omega_{jk}+ \sum^n_{l=1} p_l \delta q_l. \label{114}
\end{gather}
Therefore, the term near $\hbar^1$ becomes
\begin{equation}\label{115}\delta\left[\sum^n_{j=1}\frac{\pi {\rm i} r_j}{2}
\int^{z^{(1)}_j}_{z^{(2)}_j}\frac{\tilde{Q}}{v_0} \right]+\sum_{(jk) \in D}p_{jk} \delta \Omega_{jk}+ \sum^n_{l=1} p_l \delta q_l.\end{equation}

\textit{Coefficients near $\hbar^k$, $k \geq 2$}: by Lemma \ref{Lem_1}, we can rewrite the pairings appearing in \eqref{106} as
\begin{gather*}%\label{116}
\left\langle \oint \frac{\tilde{Q}^l}{v_0^{2l-1}}, \delta \oint \frac{\tilde{Q}^{k-l}}{v_0^{2(k-l)-1}}\right\rangle =\frac{1}{2}\int_{\partial \hat{\mathcal{C}}_0}\left( \frac{\tilde{Q}^l}{v_0^{2l-1}}
\int^x_{p_0}\delta \frac{\tilde{Q}^{k-l}}{v_0^{2(k-l)-1}}\right)+\left\langle \oint\frac{\tilde{Q}^k}{v_0^{2k-1}},\delta \oint v_0\right\rangle. \end{gather*}
Also
\begin{equation*}%\label{117}
\left\langle \oint v_0, \delta \oint \frac{\tilde{Q}^k}{v_0^{2k-1}}\right\rangle =\frac{1}{2}\int_{\partial \hat{\mathcal{C}}_0}\left(v_0
\int^x_{p_0}\delta\frac{\tilde{Q}^k}{v_0^{2k-1}}\right)+\left\langle \oint \frac{\tilde{Q}^k}{v_0^{2k-1}},\delta \oint v_0\right\rangle.
\end{equation*}
Thus, the expression \eqref{106} near $\hbar^k$, $k \geq 2$ becomes
\begin{gather}\label{118}
\binom{\frac{1}{2}}{k}\frac{1}{2}\int_{\partial \hat{\mathcal{C}}_0}\left(v_0
\int^x_{p_0}\delta\frac{\tilde{Q}^k}{v_0^{2k-1}}\right)
-\sum^{k-1}_{l=1}\binom{\frac{1}{2}}{l} \binom{\frac{1}{2}}{k-l}\frac{1}{2}\int_{\partial
\hat{\mathcal{C}}_0}\Bigg(\delta \frac{\tilde{Q}^{k-l}}{v_0^{2(k-l)-1}}
\int^x_{p_0}\frac{\tilde{Q}^l}{v_0^{2l-1}}\Bigg)\\
 \label{119}\qquad{} +\left[ \sum^{k}_{l=0}\binom{\frac{1}{2}}{l}\binom{\frac{1}{2}}{k-l}\right]
 \left\langle \oint \frac{\tilde{Q}^k}{v_0^{2k-1}},\delta \oint v_0\right\rangle .
 \end{gather}

Using the identity
\begin{equation*}%\label{120}
1+t=\big(\sqrt{1+t}\big)^2=\left(\sum^{\infty}_{k=0}\binom{\frac{1}{2}}{k}t^k \right)^2=\sum^{\infty}_{k=0}t^k \left[\sum^k_{l=0}\binom{\frac{1}{2}}{l} \binom{\frac{1}{2}}{k-l} \right], \qquad |t| \leq 1, \end{equation*}
and comparing the expressions near same powers of $t$ we conclude that the piece \eqref{119} is zero. Further, we can represent the expression in~\eqref{118} as the sum over residues at the branch points~$x_i$. Notice that the first term also has additional residues near $\big\{z^{(1)}_j, z^{(2)}_j \big\}^n_{j=1}$ due to simple poles of~$v_0$:
\begin{gather}\label{121}
\sum^n_{j=1}\binom{\frac{1}{2}}{k}\pi {\rm i}\underset{\{z^{(1)}_j, z^{(2)}_j \}}{\rm res}\left(v_0
\int^x_{p_0}\delta\frac{\tilde{Q}^k}{v_0^{2k-1}} \right)\\
 \label{122}
 {}+\sum^{4g-4+2n}_{i=1}\pi{\rm i} \, \underset{x_i}{\rm res}\left[\binom{\frac{1}{2}}{k}v_0
\int^x_{p_0}\delta\frac{\tilde{Q}^k}{v_0^{2k-1}}+ \sum^{k-1}_{l=1}\binom{\frac{1}{2}}{l} \binom{\frac{1}{2}}{k-l}\left(\frac{\tilde{Q}^l}{v_0^{2l-1}}
\int^x_{p_0}\delta \frac{\tilde{Q}^{k-l}}{v_0^{2(k-l)-1}}\right)\right].
\end{gather}
Similarly to \eqref{109}, the sum \eqref{121} could be rewritten as
\begin{equation}\label{123}\delta\left[\sum^n_{j=1}\binom{\frac{1}{2}}{k}
\pi {\rm i} r_j \int^{z^{(1)}_j}_{z^{(2)}_j}\frac{\tilde{Q}^k}{v_0^{2k-1}} \right],\end{equation}
here we used that the derivatives with respect to the coordinates $\big(A^{(0)}_i, B^{(0)}_i\big)^{g^-}_{i=1}$ commute with the integral.

Consider the expression defined on the double cover $\pi\colon \hat{\mathcal{C}}_0 \xrightarrow[]{} \mathcal{C}$, given by $v^2_0=Q_0$ in $T^*\mathcal{C}$:
\begin{equation*}%\label{124}
\underset{x_i}{\rm res}\left(\sqrt{Q_0+\hbar \tilde{Q}}\int^x_{p_0} \delta \sqrt{Q_0+\hbar \tilde{Q}}\right),
\end{equation*} where $x_i$ is a zero of $Q_0$ on~$\mathcal{C}$.

Formally, the Abelian differential $\hat{v}_\hbar=\sqrt{Q_0+\hbar \tilde{Q}}$ is globally defined on the $h$-dependent double cover $\hat{\pi}\colon \hathat{\mathcal{C}}_\hbar \xrightarrow[]{} \hat{\mathcal{C}}_0$, given by $(\hat{v}_\hbar)^2=Q_0+\hbar \tilde{Q}$ in $T^*\hat{\mathcal{C}}_0$ (note that $\hat{\mathcal{C}}_0$ itself is a double cover of $\mathcal{C}$).
Lifted from $\mathcal{C}$ to $\hat{\mathcal{C}}_0$, $Q_0$ has a 4th-order zeros at $x_i$ and double poles at $\big(z^{(1)}_j, z^{(2)}_j\big)^n_{j=1}$, while $\tilde{Q}$ gains a 2nd-order zero at $x_i$ and simple poles at $(z^{(1)}_j, z^{(2)}_j)$.
Thus, the map $\hat{\pi}$ is brached at $8g-8+4n$ simple zeroes $\tilde{x}^\hbar_j$ of ${Q_0+\hbar \tilde{Q}}$.
The double cover $\hathat{\mathcal{C}}_\hbar$ is smooth everywhere except for preimages of double zeroes $(x_i)^{4g-4+2n}_{i=1}$ of $Q_0+\hbar \tilde{Q}$, where $\hathat{\mathcal{C}}_\hbar$ gains nodes. The genus~$\hathat{g}$ of~$\hathat{\mathcal{C}}_\hbar$ equals $12g-11+4n$. Letting $\hbar \xrightarrow[]{} 0$, the nodes smoothen out and the covering surface~$\hathat{\mathcal{C}}_\hbar$ degenerates to the pair of smooth surfaces $\hat{\mathcal{C}}^{(1,2)}_0$. On the base curve $\hat{\mathcal{C}}_0$ that corresponds to the merging of triplets of points: two simple zeroes $\tilde{x}^\hbar_{i_1}$, $x^\hbar_{i_2}$ of ${Q_0+\hbar \tilde{Q}}$ converge to a double zero at~$x_i$, increasing its multiplicity to~4.

In the local coordinate $z(x)$ on $\hat{\mathcal{C}}_0$: $Q_0=v_0^2={\rm d}z^2$, $\tilde{Q}=\tilde{Q}(z){\rm d}z^2$. Differentiating with respect to the coordinates $\big(A^{(0)}_i, B^{(0)}_i\big)^{g^-}_{i=1}$ according to the rule \eqref{80}, when the coordinate $z(x)$ is kept fixed, one has that the residue could be written as
\begin{equation*}%\label{125}
\underset{x_i}{\rm res}\left(\sqrt{Q_0+\hbar \tilde{Q}}\int^x_{p_0} \frac{\hbar \delta \tilde{Q}}{2\sqrt{Q_0+\hbar \tilde{Q}}}\right)=0.
\end{equation*}
The residue vanishes since the expression inside is holomorphic at $x_i$. Then we can expand the left-hand side by powers of $\hbar$ and observe that the coefficients near the powers of $\hbar$ in the series are exactly the terms appearing in the sum \eqref{122}. It follows that these coefficients must vanish too. Thus, the coefficient near $\hbar^n, \ n \geq 2$ reduces to the expression \eqref{123}.

\textit{Full expansion}: combining \eqref{115} and \eqref{123}, we have that
\begin{gather}\label{126}\theta_1- \theta_0=\delta\left[\sum^n_{j=1}
\pi {\rm i} r_j \left[\sum^{\infty}_{k=1}\hbar^k\binom{\frac{1}{2}}{k}\int^{z^{(1)}_j}_{z^{(2)}_j}\frac{\tilde{Q}^k}{v_0^{2k-1}} \right]\right] + \hbar \left[\sum_{(jk) \in D}p_{jk} \delta \Omega_{jk}+ \sum^n_{l=1} p_l \delta q_l \right]. \end{gather}

We further notice, similarly to the argument in \eqref{101}, that the infinite series is formally the Taylor expansion by powers of $\hbar$ of the expression
\begin{equation*}%\label{127}
\int^{z^{(1)}_j(\hbar)}_{z^{(2)}_j(\hbar)}v^\hbar_1 - \int^{z^{(1)}_j}_{z^{(2)}_j}v_0 .
\end{equation*}
The issue is that the integrands are singular at the end points of the integration path, and one requires a regularization of the integral to have a proper identity.
Considering the regularization proposed in~\eqref{76}, one can see that
\begin{equation*}%\label{128}
{\rm reg}\int^{z^{(1)}_j(\hbar)}_{z^{(2)}_j(\hbar)}v^\hbar_1 ={\rm reg}\int^{z^{(1)}_j}_{z^{(2)}_j}v_0 + \sum^{\infty}_{k=1}\hbar^k \binom{\frac{1}{2}}{k}\int^{z^{(1)}_j}_{z^{(2)}_j}\frac{\tilde{Q}^k}{v_0^{2k-1}} ,
\end{equation*}
where the integrals near $h^n$ are already regular.
Using that, we rewrite the difference of potentials~as
\begin{gather}
\theta_1- \theta_0=\delta\left[\sum^n_{j=1}
\pi {\rm i} r_j \left({\rm reg}\int^{z^{(1)}_j(\hbar)}_{z^{(2)}_j(\hbar)}v^\hbar_1 - {\rm reg}\int^{z^{(1)}_j}_{z^{(2)}_j}v_0 \right)\right] \nonumber\\
\hphantom{\theta_1- \theta_0=}{}
+ \hbar \left[\sum_{(jk) \in D}p_{jk} \delta \Omega_{jk}+ \sum^n_{l=1} p_l \delta q_l \right].\label{129}
 \end{gather}

\textit{Case of two arbitrary differentials} $Q_0, Q_1 \in \mathcal{Q}_{g,n}[\mathbf{r}] $.
Theorem~1.3 of~\cite{chen2021classification} asserts that generically (outside of hyperelliptic locus for $g \geq 3$) space $\mathcal{Q}_{g,n}$, and, thus $\mathcal{Q}_{g,n}[\mathbf{r}],$ is connected. Let $\gamma_t=Q_t\colon [0,1] \xrightarrow{} \mathcal{Q}_{g,n}[\mathbf{r}] $ be a path such that $\gamma(0)=Q_0$, $\gamma(1)=Q_1.$ For each~$t$, the function $\oint_{s_i} v_t$ is holomorphic on $\mathcal{Q}_{g,n}[\mathbf{r}]$ and could be expanded by the Taylor series in some simply-connected open neighborhood $U_t$ of $Q_t$. Then $\bigcup_t U_t$ provides and open cover for $\gamma_t$. Due to compactness of $\gamma_t$, we can choose some finite subcover $\bigcup_{t_i} U_{t_i}$, $i \in [[0,N]].$ We can assume $Q_0=Q_{\hat{t}_0} \in U_{t_0}$, $Q_1=Q_{\hat{t}_{N+1}} \in U_{t_N}$ and take $Q_{\hat{t}_i} \in \gamma \cap U_{t_{i}} \cap U_{t_{i-1}}$, $i \in [[1,N]]$.
Due to~\eqref{129},
\begin{gather*}%\label{130}
\theta_{\hat{t}_{i+1}}-\theta_{\hat{t}_i}= \delta \left[\sum^n_{j=1} \pi {\rm i} r_j \left({\rm reg}\int^{z^{(1)}_j}_{z^{(2)}_j}{v_{\hat{t}_{i+1}}}-{\rm reg}\int^{z^{(1)}_j}_{z^{(2)}_j}{v_{\hat{t}_i}} \right) \right]\\
\hphantom{\theta_{\hat{t}_{i+1}}-\theta_{\hat{t}_i}=}{}
+\left[\sum_{(jk) \in D}p^{(\hat{t}_{i+1}, \hat{t}_i)}_{jk} \delta \Omega_{jk}+\sum^n_{j=1}p^{(\hat{t}_{i+1}, \hat{t}_i)}_l \delta q_l \right],
\end{gather*}
where $\big(p^{(\hat{t}_{i+1}, \hat{t}_i)}_{jk}, p^{(\hat{t}_{i+1}, \hat{t}_i)}_l\big)$ are coefficients of the linear representation of $\big(Q_{\hat{t}_{i+1}}-Q_{\hat{t}_{i}}\big)$ in the basis $\big(u_j u_k, Q^{z_l}\big)$.

Then applying the telescoping series, one has
\begin{gather}\theta_{1}-\theta_{0}=\sum^{N}_{i=0}(\theta_{\hat{t}_{i+1}}-\theta_{\hat{t}_i})= \delta \left[\sum^n_{j=1} \pi{\rm i} r_j \left({\rm reg}\int^{z^{(1)}_j}_{z^{(2)}_j}v_1-{\rm reg}\int^{z^{(1)}_j}_{z^{(2)}_j} v_0 \right) \right]\nonumber\\
\hphantom{\theta_{1}-\theta_{0}=}{}
+\left[\sum_{(jk) \in D}p^{(1,0)}_{jk} \delta \Omega_{jk}+\sum^n_{l=1}p^{(1,0)}_l \delta q_l\right],\label{131}
\end{gather}
where the latter expression, using \eqref{66}, is exactly the 1-form $\frac{1}{2}\Theta_{(S_0-S_1)}.$ By the assumption, coefficients $p^{(1,0)}_{jk}$ and $p^{(1,0)}_l$ are holomorphic functions of $(\Omega_{jk},q_l)$.
Applying differential to both sides of \eqref{131} one obtains the first statement of the proposition. If $\Theta_{(S_0 -S_1)}$ is assumed closed, then by the Poincar\'e lemma it could be locally integrated, leading to the second statement.
\end{proof}

\begin{thebibliography}{99}

\bibitem{korotkin2018periods}
Korotkin D., Periods of meromorphic quadratic differentials and {G}oldman
 bracket, in Topological Recursion and its Influence in Analysis, Geometry,
 and Topology, \textit{Proc. Sympos. Pure Math.}, Vol.~100, Amer. Math. Soc.,
 Providence, RI, 2018, 491--515, \href{https://arxiv.org/abs/1702.04705}{arXiv:1702.04705}.

\bibitem{chen2021classification}
Chen D., Gendron Q., Towards a classification of connected components of the
 strata of {$k$}-differentials, \textit{Doc. Math.} \textbf{27} (2022),
 1031--1100, \href{https://arxiv.org/abs/2101.01650}{arXiv:2101.01650}.


\end{thebibliography}
\end{document}
